\subsection{扰动}

本条中，E 是一个复希尔伯特空间。

若 \(u\) 是 E 上的自伴部分算子且 \(v \in \mathcal{L}(\mathrm{E})\) 是埃尔米特自同态，则 \(u + v\) 是自伴的（参见 IV，第 236 页）。当 \(v\) 是对称部分算子时，一般并非如此（IV，第 347 页习题 9）。然而，我们将在本条中得到这类的正面结果。

定义 11. --- 设 u 是 E 上的部分算子。如果 dom(u) ⊂ dom(v)，且 v 通过限制到该子空间定义了从 \(E_{u}\) 到 E 的连续线性映射，则称 E 上的部分算子 v 相对于 u 有界。

设 u 是 E 上的部分算子，v 是相对于 u 有界的部分算子。存在实数 m，使得

\[
\| v (x) \| \leqslant m (\| x \| + \| u (x) \|)
\]

对所有 \(x \in \text{dom}(u)\) 都成立。v 关于 u 的相对范数记为 \(\|v\|_{u}\)，它是满足以下条件的实数 \(a \geqslant 0\) 的下界：存在实数 b，使得

\[
\| v (x) \| \leqslant a \| u (x) \| + b \| x \|
\]

对所有 \(x \in \text{dom}(u)\) 都成立。因此，v 的相对范数小于 v 限制到空间 \(E_{u}\) 上所得映射的范数。

注。 --- 设 u 是 E 上的部分算子。每个自同态 \(v \in \mathcal{L}(\mathrm{E})\) 都相对于 u 有界，且其相对范数为零，因为有 \(\|v(x)\| \leqslant \|v\|\|x\|\) 对所有 \(x \in \operatorname{dom}(u)\) 都成立，这使我们可以在上面的不等式中取 a = 0。

定理 5（加藤–雷利希）。— 设 u 是 E 上的自伴部分算子，v 是相对于 u 有界的对称部分算子。若相对范数 \(\|v\|_{u}\) 为 \(< 1\)，则定义域为 dom(u) 的部分算子 \(u + v\) 是自伴的。

由于 \(\| v \|_u < 1\)，按定义存在正实数 \(a\) 和 \(b\)，使得 \(a < 1\)，且 \(\| v(x) \| \leqslant a \| u(x) \| + b \| x \|\) 对所有 \(x \in \text{dom}(u)\) 都成立。

令 \(t \in R^{*}\)。置 \(w_{t} = v \circ \mathrm{R}(u, it)\)。有 \(\operatorname{dom}(w_{t}) = E\)，因为 \(\mathrm{R}(u, it)\) 的像包含于 u 的定义域，而按假设 u 的定义域包含于 v 的定义域。设 \(x \in \operatorname{dom}(u)\)。由于 u 是自伴的，有

\[
\begin{array}{c} \| (i t - u) x \| ^ {2} = \| i t x \| ^ {2} + \| u (x) \| ^ {2} - i t \langle x | u (x) \rangle + i t \langle u (x) | x \rangle \\ = | t | ^ {2} \| x \| ^ {2} + \| u (x) \| ^ {2}, \end{array}
\]

从而有不等式 \(\|u(x)\|\leqslant\|(it-u)x\|\) 和 \(\|x\|\leqslant|t|^{-1}\|(it-u)x\|\)。于是得到

\[
\| v (x) \| \leqslant (a + b | t | ^ {- 1}) \| (i t - u) x \|.
\]

特别地，取 \(x = \mathrm{R}(u, it)y\)，其中 \(y \in \mathrm{E}\)；得到

\[
\| w _ {t} (y) \| \leqslant (a + b | t | ^ {- 1}) \| y \|.
\]

由此得到 \(w_{t}\in\mathcal{L}(\mathrm{E})\)，且 \(\|w_{t}\|\leqslant a+|t|^{-1}b\)。

由于 \(a < 1\)，存在 \(t \in \mathbf{R}_{+}^{*}\)，使得 \(a + b|t|^{-1} < 1\)。于是 \(\|w_{t}\| < 1\) 且 \(\|w_{-t}\| < 1\)；因此 E 的自同态 \(1_{\mathrm{E}} - w_{t}\) 和 \(1_{\mathrm{E}} - w_{-t}\) 都可逆（I，第 22 页命题 2）。由于有公式 \((1_{\mathrm{E}} - w_{t}) \circ (it - u) = it - (u + v)\)，这说明 \(u + v - it\) 是满射；同样，部分算子 \(u + v + it\) 是满射。因此 \(u + v\) 是自伴的（IV，第 240 页命题 11）。

